\documentclass[12pt,a4paper]{article}
\usepackage[margin=2.5cm]{geometry}
\usepackage{booktabs}
\usepackage{hyperref}
\usepackage{amsmath}
\usepackage{xcolor}
\usepackage{listings}
\usepackage{microtype}
\usepackage{parskip}

\lstset{
  basicstyle=\ttfamily\small,
  breaklines=true,
  frame=single,
  backgroundcolor=\color{gray!8},
}

\definecolor{warn}{RGB}{180,0,0}
\newcommand{\warn}[1]{\textcolor{warn}{\textbf{[!] #1}}}

\title{\textbf{Cost Model: Codebase Audit}\\[0.4em]
  \large \texttt{netsquid-qkd} --- 2026-08-03}
\author{}
\date{}

\begin{document}
\maketitle
\tableofcontents
\newpage

%==============================================================================
\section{Overview}
%==============================================================================

All monetary values are in GBP. The cost model has two components:

\begin{itemize}
  \item \texttt{network/cost.py} --- core module defining unit costs, technology
        presets, component count formulae, the SPD linear cost model, and the
        \texttt{total\_cost} function.
  \item \texttt{scripts/network/cost\_sweep.py} --- analysis script that sweeps
        user count $N$ at fixed relay count $K$, runs simulations for all three
        protocols, and plots deployment cost and cost-efficiency.
\end{itemize}

The SQLite database (\texttt{data/results.db}) stores fibre length from each
simulation run but no cost figures. Cost is recomputed on demand by
\texttt{cost\_sweep.py} using the parameters passed at the time of analysis.

%==============================================================================
\section{Logic Flow}
%==============================================================================

\subsection{File dependency graph}

\begin{verbatim}
network/cost.py
  +-- defines: DEFAULT_COSTS, DETECTOR_TECH, SOURCE_TECH
               spd_cost_from_efficiency(eta)
               component_counts(N, K, protocol)
               total_cost(counts, total_fibre_km, detector_efficiency, ...)

network/bb84_network.py
  +-- run_bb84_network()   -> returns dict including total_fibre_km
network/mdi_network.py
  +-- run_mdi_network()    -> returns dict including total_fibre_km
network/trusted_bb84_network.py
  +-- run_trusted_bb84_network() -> returns dict including total_fibre_km

lib/db.py
  +-- network_results table: stores total_fibre_km (no cost columns)
  +-- p2p_results table:     stores per-pair physics results (no cost)

scripts/network/cost_sweep.py
  +-- imports cost.py; calls run_*_network(); calls component_counts()
      + total_cost() to compute cost; plots cost vs N, efficiency vs N,
      and marginal cost vs N
\end{verbatim}

\subsection{Cost computation pattern in \texttt{cost\_sweep.py}}

\begin{enumerate}
  \item Run QKD simulation to obtain \texttt{res["total\_fibre\_km"]}.
  \item Read \texttt{cfg["detector\_efficiency"]} (set by \texttt{--detector-tech}
        preset or simulation config).
  \item Call \texttt{component\_counts(N, K, protocol)} for hardware quantities.
  \item Call \texttt{total\_cost(counts, fibre\_km, det\_eff, **cost\_kw)}.
  \item Read \texttt{["total\_gbp"]} from the result.
\end{enumerate}

In \texttt{--cost-only} mode the simulation is skipped; fibre length is computed
directly from Euclidean topology geometry via \texttt{\_fibre\_km\_bb84(topo)}
and \texttt{\_fibre\_km\_mdi(topo)}.

\subsection{Fibre distance calculation per protocol}

\begin{description}
  \item[BB84 direct mesh] Sum of all $N(N{-}1)/2$ pairwise Euclidean distances.
        Models a fully meshed dark fibre deployment with one fibre per user pair.
  \item[MDI-QKD] $N$ user--relay links (Euclidean) plus $K(K{-}1)/2$
        relay--relay backbone links (Euclidean). Computed by
        \texttt{\_mdi\_fibre\_cost()} in \texttt{mdi\_network.py}.
  \item[Trusted BB84] Identical topology to MDI: $N$ user--relay links plus
        $K(K{-}1)/2$ relay--relay backbone links. The backbone carries
        quantum signals (relay-to-relay BB84) as well as classical key
        forwarding traffic.
\end{description}

All distances are straight-line Euclidean on a flat $A \times A$\,km grid.
No tortuosity or routing factor is applied.

%==============================================================================
\section{SPD Cost Model}
%==============================================================================

\subsection{Linear model}

Rather than assigning a fixed price per detector, SPD cost is modelled as a
linear function of detector efficiency $\eta$, anchored at two real technology
classes:

\begin{table}[h!]
\centering
\begin{tabular}{llcc}
\toprule
Key & Technology & Efficiency $\eta$ & Cost per SPD (GBP) \\
\midrule
\texttt{SPAD}  & InGaAs SPAD, telecom-band (1550\,nm) & 0.20 & \pounds15{,}000 \\
\texttt{SNSPD} & Superconducting nanowire, cryogenic ($\sim$2\,K) & 0.85 & \pounds100{,}000 \\
\bottomrule
\end{tabular}
\caption{\texttt{DETECTOR\_TECH} anchor points in \texttt{network/cost.py}.}
\end{table}

The linear fit through these two points gives:

\[
C_{\mathrm{SPD}}(\eta) = m\,\eta + c
\]

where

\[
m = \frac{100{,}000 - 15{,}000}{0.85 - 0.20} = \frac{85{,}000}{0.65}
  \approx \pounds130{,}769 \text{ per unit }\eta
\]
\[
c = 15{,}000 - m \times 0.20 \approx -\pounds11{,}154
\]

This is implemented in \texttt{spd\_cost\_from\_efficiency(eta)}, which raises a
\texttt{ValueError} for $\eta < 0.20$ (below the InGaAs SPAD baseline).
The function is called internally by \texttt{total\_cost}; there is no
explicit per-SPD cost override available at the CLI level.

Spot checks: $\eta = 0.20 \to \pounds15{,}000$;\quad
$\eta = 0.85 \to \pounds100{,}000$;\quad
$\eta = 0.50 \to \pounds54{,}231$.

%==============================================================================
\section{Component Model}
%==============================================================================

\subsection{Hardware quantities per protocol}

The function \texttt{component\_counts(N, K, protocol)} returns the following:

\begin{table}[h!]
\centering
\begin{tabular}{lccccccl}
\toprule
Protocol & Sources & SPDs & BS & PBS & EOM & Switch & Links \\
\midrule
BB84 (direct mesh) & $N$ & $2N$ & $0$ & $N$ & $N$ & $0$ & $N(N{-}1)/2$ \\
MDI-QKD            & $N$ & $4K$ & $K$ & $2K$ & $0$ & $K$ & $N + K(K{-}1)/2$ \\
Trusted BB84       & $N+K$ & $2K$ & $0$ & $K$ & $K$ & $K$ & $N + K(K{-}1)/2$ \\
\bottomrule
\end{tabular}
\caption{Hardware component counts. $N$ = users, $K$ = relays.
BS = 50:50 fibre coupler; PBS = polarising beam splitter;
EOM = electro-optic modulator.}
\end{table}

\subsection{Physical rationale per protocol}

\paragraph{BB84 direct mesh.}
Each user has one photon source (Alice) and a receiving module (Bob).
Active basis choice is assumed: an EOM rotates the incoming polarisation
to align with a fixed PBS, effectively switching which basis is measured.
Bob components per user: 1 EOM + 1 PBS + 2 SPDs. Full mesh: one dark
fibre per pair, hence $N(N{-}1)/2$ links.

\paragraph{MDI-QKD.}
Users supply sources; MDI relay nodes are passive Bell state measurement
(BSM) stations --- no source at a relay, no basis choice required.
Each relay performs a full polarisation-resolving BSM:
\begin{itemize}
  \item 1 $\times$ 50:50 fibre coupler --- Hong--Ou--Mandel interference.
  \item 2 $\times$ PBS --- one per output arm, for polarisation analysis.
  \item 4 $\times$ SPDs --- two per PBS (H and V outputs), giving complete
        Bell state discrimination between $|\Psi^+\rangle$ and $|\Psi^-\rangle$.
  \item 1 $\times$ optical switch --- routes photons from cross-cluster users
        to the correct relay for BSM.
\end{itemize}
Fibre links: $N$ user-to-relay links plus $K(K{-}1)/2$ relay-to-relay backbone links.

\paragraph{Trusted BB84.}
Users connect to their assigned relay via direct BB84 links. The relay acts
as Bob on all user-relay links, sharing a single detector array (1 EOM +
1 PBS + 2 SPDs) across all connected users. An optical switch routes each
incoming user signal to the detector array in turn. Relays also run BB84
with one another over the backbone to establish inter-relay keys, requiring
one photon source per relay (simplified from the strict $K{-}1$ sources per
relay that a fully connected backbone would require). Total sources: $N + K$.
No BSM hardware. Classical key forwarding (XOR masking) is authenticated
but not costed (see Section~\ref{sec:classical}).

\subsection{Trusted BB84 relay-relay key security note}

The classical message transmitted between relay nodes during key forwarding
takes the form $K_{AR_1} \oplus K_{R_1 R_2}$, where $K_{R_1 R_2}$ was
established by relay-to-relay QKD and is secret. This is equivalent to a
one-time pad: an in-flight attacker observing the classical channel cannot
recover $K_{AR_1}$. The true trust assumption is not channel security
but node security: a physically compromised relay node has access to all
key material it has processed. MDI-QKD eliminates this vulnerability because
relay nodes (Charlie) never hold any key material.

%==============================================================================
\section{Total Cost Formula}
%==============================================================================

\[
C_{\mathrm{total}} = N_s C_s
                   + N_d\,C_{\mathrm{SPD}}(\eta)
                   + N_b C_b
                   + N_{\mathrm{pbs}} C_{\mathrm{pbs}}
                   + N_{\mathrm{eom}} C_{\mathrm{eom}}
                   + N_{\mathrm{sw}} C_{\mathrm{sw}}
                   + L_{\mathrm{fibre}} C_f
\]

where $N_s, N_d, N_b, N_{\mathrm{pbs}}, N_{\mathrm{eom}}, N_{\mathrm{sw}}$
are the source, SPD, 50:50 BS, PBS, EOM, and optical switch counts from
\texttt{component\_counts}; $C_{\mathrm{SPD}}(\eta)$ is the linear SPD cost
model; $L_{\mathrm{fibre}}$ is total fibre in km; and remaining $C_x$ are
unit costs.

\texttt{total\_cost()} returns a dict with \texttt{total\_gbp},
\texttt{hardware\_gbp}, \texttt{fibre\_gbp}, and a \texttt{breakdown} sub-dict
with per-category figures.

%==============================================================================
\section{Default Costs and Technology Presets}
%==============================================================================

\subsection{\texttt{DEFAULT\_COSTS}}

\begin{table}[h!]
\centering
\begin{tabular}{llrl}
\toprule
Key & Component & Default (GBP) & Notes \\
\midrule
\texttt{source\_gbp}         & Photon source (QD)  & \pounds150{,}000  & QD SPS + cryostat, amortised \\
\texttt{bs\_gbp}             & 50:50 fibre coupler & \pounds1{,}000    & HOM BS at MDI relay \\
\texttt{pbs\_gbp}            & Polarising BS       & \pounds500        & BB84/TBB84 receivers; MDI arms \\
\texttt{eom\_gbp}            & EOM                 & \pounds2{,}000    & Active basis choice; BB84/TBB84 \\
\texttt{switch\_gbp}         & Optical switch      & \pounds5{,}000    & Cross-relay routing; MDI only \\
\texttt{fibre\_per\_km\_gbp} & Installed fibre     & \pounds10{,}000/km & Dark fibre, installed \\
\bottomrule
\end{tabular}
\caption{\texttt{DEFAULT\_COSTS} in \texttt{network/cost.py}. SPD cost is not
listed here --- it is always derived from detector efficiency via the linear
model. All non-SPD values are indicative placeholders pending cost verification.}
\end{table}


%==============================================================================
\section{CLI Flags in \texttt{cost\_sweep.py}}
%==============================================================================

\begin{table}[h!]
\centering
\begin{tabular}{lll}
\toprule
Flag & Type & Effect \\
\midrule
\texttt{--detector-tech SPAD|SNSPD} & preset & Sets $\eta$ in sim; SPD cost derived from $\eta$ \\
\texttt{--source-cost FLOAT} & GBP override & Replaces QD default source cost \\
\texttt{--bs-cost FLOAT}     & GBP override & Replaces default 50:50 BS cost \\
\texttt{--pbs-cost FLOAT}    & GBP override & Replaces default PBS cost \\
\texttt{--eom-cost FLOAT}    & GBP override & Replaces default EOM cost \\
\texttt{--switch-cost FLOAT} & GBP override & Replaces default optical switch cost \\
\texttt{--fibre-cost FLOAT}  & GBP/km override & Replaces default fibre cost \\
\bottomrule
\end{tabular}
\caption{Cost-relevant CLI flags. There is no \texttt{--spd-cost} or
\texttt{--source-tech} flag; SPD cost is always derived from detector
efficiency, and source technology is fixed to QD.}
\end{table}

Priority order: \texttt{DEFAULT\_COSTS} $\to$ \texttt{--detector-tech} preset
$\to$ explicit CLI flag. Detector efficiency is taken from
\texttt{cfg["detector\_efficiency"]}, set by \texttt{--detector-tech} or the
JSON config file.

%==============================================================================
\section{Database Schema and Cost Persistence}
%==============================================================================

\subsection{What is stored}

\textbf{\texttt{network\_results}} --- one row per network simulation run:
\begin{itemize}
  \item Protocol, $N$, $K$, area, seed, runtimes, experiment label.
  \item \texttt{total\_fibre\_km}, \texttt{avg\_pair\_distance\_km},
        \texttt{cross\_relay\_ratio}.
  \item Key rate statistics: \texttt{avg\_key\_rate}, \texttt{std\_key\_rate},
        \texttt{success\_rate}, \texttt{min\_key\_rate}, \texttt{max\_key\_rate}.
\end{itemize}

\textbf{\texttt{p2p\_results}} --- one row per Monte Carlo runtime per pair:
physical parameters, key rate, QBER, key length, status.

\subsection{What is not stored}

No cost column exists in either table. The unit costs used at analysis time
(source, SPD efficiency, BS, PBS, switch, fibre) are not persisted alongside
simulation results. A stored network row cannot be repriced from the database
alone without knowing the cost parameters supplied at the time.

\textbf{Proposed addition} --- add the following columns to
\texttt{network\_results} at insert time:

\begin{lstlisting}[language=SQL]
ALTER TABLE network_results ADD COLUMN source_tech          TEXT;
ALTER TABLE network_results ADD COLUMN detector_tech        TEXT;
ALTER TABLE network_results ADD COLUMN detector_efficiency  REAL;
ALTER TABLE network_results ADD COLUMN source_cost_gbp      REAL;
ALTER TABLE network_results ADD COLUMN spd_cost_gbp         REAL;
ALTER TABLE network_results ADD COLUMN bs_cost_gbp          REAL;
ALTER TABLE network_results ADD COLUMN pbs_cost_gbp         REAL;
ALTER TABLE network_results ADD COLUMN switch_cost_gbp      REAL;
ALTER TABLE network_results ADD COLUMN fibre_cost_per_km_gbp REAL;
ALTER TABLE network_results ADD COLUMN hardware_cost_gbp    REAL;
ALTER TABLE network_results ADD COLUMN fibre_cost_gbp       REAL;
ALTER TABLE network_results ADD COLUMN total_cost_gbp       REAL;
\end{lstlisting}

%==============================================================================
\section{Assumptions and Limitations}
%==============================================================================

\label{sec:classical}
\subsection{Classical communication is free}

All classical communication costs --- synchronisation channels, basis
reconciliation, error correction, privacy amplification, authenticated key
forwarding at trusted relays --- are assumed to be zero. Only quantum-layer
hardware and quantum fibre links are costed. This is a standard simplification
for quantum network techno-economic analysis at the early-stage modelling level.

\subsection{Fibre is Euclidean, not routed}

All link distances are straight-line Euclidean. Real fibre follows duct
paths; a tortuosity factor of $1.2$--$1.5\times$ is typical in urban
deployments. A \texttt{--tortuosity} flag (default 1.0) could scale all
fibre distances before cost computation without affecting the physics simulation.

\subsection{BB84 mesh fibre}

The BB84 cost model counts $N(N{-}1)/2$ individual dark fibre runs, one per
user pair. This is correct for a fully meshed passive dark-fibre network but
unrealistic at scale; a WDM or switched architecture would replace individual
fibres. This makes BB84 costs appear very large at high $N$, which is physically
meaningful (mesh BB84 does not scale well) but should be interpreted with caution
as a direct deployment cost estimate.

\subsection{Cost and simulation models are not fully coupled}

The cost model counts individual optical components (EOM, PBS, 50:50 BS,
optical switch) at each node. The simulation, however, abstracts all
receiver-side optics into a single \texttt{node\_loss\_db} parameter and a
single \texttt{detector\_efficiency} value. Loss contributions from the EOM,
PBS, and switching fabric are not modelled separately in the physics layer ---
they are subsumed into the aggregate node loss figure. Similarly, the optical
switch at trusted BB84 and MDI relays introduces insertion loss in practice,
but this is not broken out as a distinct simulation parameter.

This means the cost model and the rate model are not tightly coupled: adding
more components raises the modelled cost but does not automatically degrade
the simulated key rate. The correct approach when refining component budgets
is to update both the cost model here \emph{and} the node loss parameter in
the simulation config to reflect realistic insertion losses for the chosen
hardware.

\subsection{No operational expenditure}

The model is purely capital expenditure (CapEx). For cryogenic technologies
(SNSPD at $\sim$2\,K, QD at $\sim$4\,K) the operational cost of cryogenic
cooling can rival the hardware CapEx over a 10-year lifetime at commercial
electricity prices, and is entirely unmodelled here.

\subsection{Source technology is fixed to QD}

The simulation uses ideal single-photon sources throughout. For cost purposes,
the source is assumed to be a semiconductor quantum dot (QD) system at
\pounds150{,}000 per node, reflecting cryogenic hardware and amortised cryostat
cost. This value is a literature-derived estimate and should be verified.
Use \texttt{--source-cost} to override when a revised figure is available.

\subsection{Source count is approximate for trusted BB84 backbone}

The model assigns $K$ backbone sources (one per relay) for relay-to-relay
QKD links. Strictly, a fully connected backbone with $K$ nodes and
bidirectional BB84 requires each relay to have $K{-}1$ sources (one per
neighbour), giving $K(K{-}1)$ total backbone sources. The current approximation
of 1 source per relay underestimates backbone hardware cost for $K > 2$.
For the relay counts typical in this project ($K \leq 8$) the error is
bounded but non-negligible at high $K$.

%==============================================================================
\section{Cross-Analysis: Simulation vs.\ Cost Model}
%==============================================================================

This section maps every simulation parameter in \texttt{lib/functions.py}
and the network runner files against its counterpart (or absence) in the cost
model, and documents all identified inconsistencies.

\subsection{Simulation parameters and their cost model counterparts}

\begin{table}[h!]
\centering
\small
\begin{tabular}{p{4.2cm}p{3.0cm}p{2.5cm}p{4.2cm}}
\toprule
Simulation parameter & Default value & Cost model entry & Notes \\
\midrule
\texttt{fibre\_loss\_db\_per\_km} & 0.20 dB/km & \texttt{fibre\_per\_km\_gbp} & Physics and cost are independent; coupled only by total km \\
\texttt{init\_loss} & 0.15 (linear) & \emph{none} & TX-side coupling loss; no cost component \\
\texttt{detector\_efficiency} & 0.85 & \texttt{spd\_cost\_from\_efficiency} & Fully coupled: $\eta$ drives both rate and SPD cost \\
\texttt{dark\_count\_rate} & 100 cps & \emph{none} & Only $\eta$ drives cost; dark count not linked to detector class \\
\texttt{node\_loss\_db} & 2.0 dB & EOM + PBS + switch + BS & All RX-side components lumped into one scalar \\
\texttt{source\_error\_rate} & 0.005 & \emph{none} & Not linked to QD assumption \\
\texttt{dephasing\_rate} & 0.0001 & \emph{none} & Physical noise; no hardware cost correlate \\
\texttt{bs\_eff} & 0.97 (MDI only) & \texttt{bs\_gbp} & Physics and cost independent; both MDI-only \\
\texttt{SWITCH\_LOSS\_DB} & 1.0 dB (MDI) & \texttt{switch\_gbp} & Loss conditional; cost unconditional (see below) \\
\bottomrule
\end{tabular}
\caption{All simulation parameters against cost model entries. \emph{None}
indicates no cost model counterpart exists.}
\end{table}

\subsection{Inconsistency I: \texttt{node\_loss\_db} aggregates individually-costed components}

\textbf{Severity: high.}

The simulation aggregates all receiver-side optical losses into a single
\texttt{node\_loss\_db} parameter (default 2.0\,dB, described in
\texttt{configs/layer5\_node\_loss.json} as ``FC/PC connectors + optical
components at receiver''). The cost model, by contrast, prices each component
individually: EOM, PBS, 50:50 BS (MDI), and optical switch.

Typical per-component insertion losses are:
\begin{itemize}
  \item EOM (lithium niobate): 0.5--3.0\,dB
  \item PBS (fibre-coupled): 0.1--0.5\,dB
  \item Optical switch (MEMS or thermo-optic): 0.5--2.0\,dB
  \item 50:50 fibre coupler (HOM BS): 0.1--0.3\,dB (excess loss only; 3\,dB splitting is accounted for by \texttt{bs\_eff})
\end{itemize}

For a BB84 or trusted BB84 receiver with EOM + PBS, the aggregate RX loss
could plausibly reach 3.5--3.5\,dB. For an MDI relay with 50:50 BS + PBS +
switch, aggregate loss could reach 2.6--4.3\,dB. The 2.0\,dB default is
therefore potentially too low, and adding components to the cost model
without updating \texttt{node\_loss\_db} in the simulation configs will
silently overstate performance while raising cost.

\textbf{Recommended action:} define per-component loss budgets and update
\texttt{node\_loss\_db} in the relevant JSON configs to reflect the actual
hardware stack being costed.

\subsection{Inconsistency II: SPAD dark count rate not updated with detector class}

\textbf{Severity: high.}

When \texttt{--detector-tech SPAD} is passed, the simulation correctly sets
$\eta = 0.20$, but \texttt{dark\_count\_rate} remains at its default of
100\,cps. This is unrealistic: InGaAs SPADs have dark count rates of
approximately 1{,}000--50{,}000\,cps, compared to 1--100\,cps for SNSPDs.
Simulating SPAD cost with SNSPD-tier dark counts significantly overstates
SPAD key rate performance.

\textbf{Recommended action:} couple \texttt{dark\_count\_rate} to
\texttt{--detector-tech} in \texttt{cost\_sweep.py}, e.g.\
SPAD $\to$ 10{,}000\,cps, SNSPD $\to$ 100\,cps. These values should
themselves be verified against the literature.

\subsection{Inconsistency III: source error rate not linked to QD physics}

\textbf{Severity: low--medium.}

The simulation default \texttt{source\_error\_rate = 0.005} (0.5\%) is a
fixed value with no explicit connection to the QD source assumption used in
the cost model. QD single-photon sources have finite multi-photon emission
probability ($g^{(2)}(0) > 0$) that contributes to QBER. A value of 0.5\%
is plausible for a high-quality QD device but is not derived from a specific
QD literature value. There is currently no mechanism to propagate source
technology choice into the simulation physics.

\textbf{Recommended action:} note the assumed value in the simulation config
or docstring, and cite the QD literature source when verifying the cost figure.

\subsection{Inconsistency IV: stale \texttt{n\_spd} in MDI network return dict}

\textbf{Severity: low (documentation only).}

\texttt{mdi\_network.run\_mdi\_network()} returns a dict that includes
\texttt{"n\_spd": 2 * topo.K}. The cost model now uses $4K$ SPDs per MDI
deployment (full polarisation-resolving BSM). The return dict value is not
used in cost computation --- cost is computed via \texttt{component\_counts}
independently --- but the stale field is misleading if the return dict is
used for diagnostics or future analysis.

\textbf{Recommended action:} update the return dict to \texttt{"n\_spd": 4 * topo.K}.

\subsection{Inconsistency V: optical switch loss is pair-conditional; cost is per-relay}

\textbf{Severity: low (by design, but worth noting).}

In \texttt{mdi\_network.py}, the optical switch insertion loss
(\texttt{SWITCH\_LOSS\_DB = 1.0\,dB}) is applied only to cross-cluster pairs,
i.e.\ pairs whose two users are assigned to different relays. In-cluster pairs
(both users on the same relay) incur no switch loss. The cost model, however,
charges for one switch per MDI relay unconditionally.

This is physically correct --- the switch must be installed at every relay
regardless of which pairs use it --- but it creates an asymmetry: the cost
is borne by the network as a whole while the performance penalty is borne only
by cross-cluster pairs. Networks with high $K$ (many relays, low cross-cluster
ratio) will show a larger switch cost fraction relative to the switch-induced
rate penalty.

\subsection{Inconsistency VI: TX-side coupling loss has no cost entry}

\textbf{Severity: negligible.}

\texttt{init\_loss = 0.15} models source-to-fibre coupling efficiency
(15\% loss). The coupling hardware --- fibre collimator, circulator, or
pigtailed connector at the transmitter --- has no corresponding entry in
the cost model. Unit costs for such components are typically
\pounds100--\pounds500 each, small relative to the QD source at \pounds150k.
Omitting them is acceptable but should be acknowledged.

\subsection{Inconsistency VII: trusted BB84 backbone is simulated as unidirectional}

\textbf{Severity: medium.}

\texttt{trusted\_bb84\_network.py} simulates $K(K{-}1)/2$ backbone links,
one per relay pair, each running a single unidirectional BB84 instance.
The key generated on link $(j, k)$ is shared between relay $j$ and relay $k$,
and is divided among all cross-relay end-to-end pairs that traverse that
backbone segment (key dilution). For full bidirectional key generation ---
where $j$ can independently forward a key to $k$ and $k$ can forward a
separate key to $j$ --- each relay would need two BB84 sessions per neighbour
(doubling backbone source count and rate). The simulation implicitly treats
the single key as shared, which may overstate backbone throughput for dense
cross-relay traffic.

\subsection{Inconsistency VIII: dark count rate absent from SPD cost model}

\textbf{Severity: low.}

The linear SPD cost model $C_{\mathrm{SPD}}(\eta) = m\eta + c$ depends only
on detector efficiency. Dark count rate is a second major performance
parameter that also varies strongly with technology class: SNSPDs achieve
$\sim$1--100\,cps while InGaAs SPADs are typically $\sim$1{,}000--50{,}000\,cps.
Both efficiency and dark count rate are correlated with cost in practice, but
the model captures only the efficiency dimension. This means two detector
technologies with the same $\eta$ but different dark count rates would be
assigned identical costs, which may not reflect reality for technologies
at intermediate efficiencies on the linear model.

%==============================================================================
\section{Quick Reference: Effective Costs at Key Configurations}
%==============================================================================

\subsection{SPAD detectors, QD source}

\begin{table}[h!]
\centering
\begin{tabular}{lrl}
\toprule
Component & Cost (GBP) & Origin \\
\midrule
Source (QD, per node)          & \pounds150{,}000 & \texttt{DEFAULT\_COSTS} \\
SPD (SPAD, per detector)       & \pounds15{,}000  & \texttt{spd\_cost\_from\_efficiency(0.20)} \\
EOM (BB84/TBB84 receivers)     & \pounds2{,}000   & \texttt{DEFAULT\_COSTS} \\
PBS                            & \pounds500       & \texttt{DEFAULT\_COSTS} \\
50:50 BS (MDI relay)           & \pounds1{,}000   & \texttt{DEFAULT\_COSTS} \\
Optical switch (MDI relay)     & \pounds5{,}000   & \texttt{DEFAULT\_COSTS} \\
Fibre (per km, installed)      & \pounds10{,}000  & \texttt{DEFAULT\_COSTS} \\
Detector efficiency $\eta$     & 0.20             & \texttt{DETECTOR\_TECH["SPAD"]} \\
\bottomrule
\end{tabular}
\caption{Run with \texttt{--detector-tech SPAD}. Optical switch applies to both
MDI and trusted BB84 relays.}
\end{table}

\subsection{SNSPD detectors, QD source}

\begin{table}[h!]
\centering
\begin{tabular}{lrl}
\toprule
Component & Cost (GBP) & Origin \\
\midrule
Source (QD, per node)          & \pounds150{,}000 & \texttt{DEFAULT\_COSTS} \\
SPD (SNSPD, per detector)      & \pounds100{,}000 & \texttt{spd\_cost\_from\_efficiency(0.85)} \\
EOM (BB84/TBB84 receivers)     & \pounds2{,}000   & \texttt{DEFAULT\_COSTS} \\
PBS                            & \pounds500       & \texttt{DEFAULT\_COSTS} \\
50:50 BS (MDI relay)           & \pounds1{,}000   & \texttt{DEFAULT\_COSTS} \\
Optical switch (MDI relay)     & \pounds5{,}000   & \texttt{DEFAULT\_COSTS} \\
Fibre (per km, installed)      & \pounds10{,}000  & \texttt{DEFAULT\_COSTS} \\
Detector efficiency $\eta$     & 0.85             & \texttt{DETECTOR\_TECH["SNSPD"]} \\
\bottomrule
\end{tabular}
\caption{Run with \texttt{--detector-tech SNSPD}.}
\end{table}

\end{document}
